\section{Who May Amend the Floor?}
\label{sec:amend-floor}
\runin{Where an error goes}
The floor's structural position is the unpublished administrative
disposition rather than the judge's: unreasoned, unindexed,
non-precedential, resolved in volumes no forum sees, where identical
errors accumulate without becoming a pattern anybody can point at,
the affected party has no way to learn that others are in the same
position, and the institution is the only repeat player. The
comparison with a constitutional court shows why that is worse than
it sounds. A Karlsruhe error is national, immediate and uniform: the
Court's decisions bind every court and every authority of the
Federation and the Länder, and in the review proceedings carry the
force of law \autocite{germany1951bverfgg}. But binding force is what
makes an error correctable, because it makes the error recur;
recurrence produces parties with a reason to litigate it, and
recurrence across a boundary produces a disagreement between courts
that a higher court has to settle. The American structure corrects by
percolation; the Federal Republic, lacking the boundary, corrects by
amendment, and amends its constitution often. A floor's error is as
uniform as Karlsruhe's and recurs as reliably, and nothing converts
the recurrence into visibility. A second deployment holding the same
floor adds affected parties and no forum. The floor chapter builds a
route for one such party. It does not build the thing that would let
a thousand of them find each other.

The other difference is custody rather than propagation. The Court
did not write the statute it reads; it can be overridden by a
supermajority that has been assembled sixty-odd times since 1949; the
state it belongs to is bound by the court at Strasbourg and by Union
law as the court at Luxembourg reads it; and its two Senates must go
to a plenum when they diverge. An operator writes the floor, applies
it, and answers to nobody. That is the capturable setting
the previous section named.

\runin{What amending it would take}
The proposed custody measures — publication, threshold control of retraining, attestation, an adverse-interest quorum, and a removal bond — can also govern amendment. Nothing distinguishes an amendment from a removal at the level of what happens to the artifact; both end with a model that no longer refuses what it refused before. What distinguishes them is how many parties had to agree, whether the attempt is on the record when it fails, and whether the change was declared before it operated. So the same five terms are the procedure, and what was missing was never the mechanism. It was the decision to run it in the open.

Three conditions follow and two are already argued for. The number of parties who must agree to a change has to exceed the number needed to run the deployment. The attempt has to reach the record whether it carries or fails, and a proposal that fails moves no weights, so the parties who had to agree are the ones who record that they were asked. The failed attempt is the more informative of the two: a floor nobody has tried to change and a floor whose removal was proposed and refused look identical from outside, and only the record separates them. The third condition is new here. An amended floor binds forward and does not reach back: a model retrained to permit an act cannot then be offered as evidence that the act was never covered, which it otherwise can be, the artifact that would contradict it being the one that was replaced. Non-retroactivity protects the record and empties the correction. It is what keeps an operator from rewriting what the floor said, and it is also why nobody harmed while the wrong version was operating gets anything from the amendment that fixed it. Adjudication has the same pair under the name of prospective overruling: the corrected rule binds forward, and the parties decided under the old one keep what they got. The condition is therefore a trade rather than a safeguard, and a regime that wanted the second half would need a remedy running backward, which nothing here supplies. Versioning the published list and dating the attestation is what makes that checkable, and it asks nobody to do anything they were not doing already.

The nearest instrument outside this book sits in the tradition named above. The Basic Law the Federal Republic adopted in 1949 requires two thirds of both chambers for any amendment at all, and then places a short list of principles beyond amendment entirely \autocite{germany1949basiclaw}. It was drafted against a particular memory, the 1933 statute by which a legislature voted away its own power to check what came next, and it is militant democracy given a text rather than an argument. A floor can have both layers in the same way: a published list that changes under a procedure, and a core the procedure cannot reach. What it cannot have is an account of who decides which is which that does not run straight back into this section's problem.

One disanalogy makes the floor's version harder than the constitutional one, and it is the same fact about deployment that makes a single bearer's error uniform. An amendment to a constitution reaches every future case because every court reads the same text. An amendment reaches a deployment only if it is updated, so the procedure necessarily retains a channel the operator could use for silent revision. A procedure that shut the channel would be one under which no deployment could be corrected at all. So no version of this procedure closes it, and the exposure is the price of every correction the procedure exists to make.

The amendment procedure is not an appeal route: its participants already hold authority over the floor. The wrongful-refusal procedure addresses complaints from those harmed by a refusal. The party the floor claims to protect reaches neither, and the two things it cannot do want different answers. Declining the protection is a waiver obtained from a party whose circumstances somebody else arranged, a floor the protected party can decline is a floor an operator can arrange to have declined. Initiating review of the contents is the half that stays open, and the difficulty is not in drafting a procedure. A route by which the protected population sets what the floor contains makes the floor an aggregation output, and independence from the aggregate is what it was for.

So the claim of authority is not discharged, and what this section and the last license is narrower than a floor. Everything it endorses can be checked from outside the operator, and being checkable is not being entitled. What it does not license is the occasion the reply above runs on: an operator that has found for itself that the emergency obtains has not thereby acquired the authority to install a floor.

Entitlement of that kind is not acquired at the moment of writing, and the American record has one instance of how it is acquired instead. The most consequential floor in the American record was published by men who had just struck from it the clause that would have bound them, at two delegations' insistence and with the reason recorded by its drafter: the passage reprobating the slave trade went, Jefferson wrote, “in complaisance to South Carolina \& Georgia,” with the northern delegates who carried the trade feeling tender under the censure too \autocite{loc2010declaration}. What remained was general enough to be turned on them. It was being turned on them already — Black petitioners in Massachusetts had been making the natural-rights argument since 1773, and within six months of the Declaration eight of them were making it in the Declaration's own terms, claiming “in common with all other Men, a natural and unalienable right to that freedom” \autocite{hall1777petition}. The legislature did not act. What eventually made the proposition binding was not that its authors had found it self-evident, which their own deletion shows they had not, but that the omission was on the record, and others kept holding the remainder to them at a price. Held to it, in the end, by a war and three amendments. What made the proposition binding was enforcement; what the record supplied was the object the enforcement could be aimed at, and a record no lever ever reached is a petition the legislature did not act on.

A floor is in that position. The parties who will hold its author to the remainder are not the people it protects, who have no route by which to do it; they are whoever else writes a floor. A published floor is legible to the next party drafting, an omission is legible beside it, and an operator that holds a term when it bites sets a price the next operator is measured against. A floor is written by anyone who installs a constraint, and continuous adaptation multiplies the parties who install them. That constituency is not the protected population, and the closing chapter says what reviewing oneself costs; it is nonetheless the only one positioned to hold this floor to anything.
